%% -----------------------------------------------------------------
%% 北京邮电大学研究生学位论文 - 第三章 实验与结果分析
%% -----------------------------------------------------------------

\chapter{实验验证与结果分析}

\section{实验环境与数据集}

为了验证本文所提理论模型与算法方案的优越性，本章在仿真测试平台与实际工程语料上进行了充分的实验评估。基准测试数据集包含来自真实通信链路的连续采样样本。

\begin{table}[htbp]
  \centering
  \bicaption{不同算法在测试集上的性能对比指标}{Performance Comparison of Different Algorithms on the Benchmark Dataset}
  \label{tab:exp_results}
  \renewcommand{\arraystretch}{1.3}
  \begin{tabular}{lcccc}
    \toprule
    算法方案 & 运行时间 ($\mathrm{ms}$) & 吞吐量 ($\mathrm{Mbps}$) & 准确率 (\%) & 能效比 ($\mathrm{bit/J}$) \\
    \midrule
    基线算法 A~\cite{nalamura2014} & 45.2 & 124.6 & 88.5 & 312.4 \\
    基线算法 B~\cite{mahuan2011}   & 38.7 & 142.1 & 91.2 & 389.0 \\
    本文所提方案                    & \textbf{21.4} & \textbf{198.5} & \textbf{96.8} & \textbf{542.7} \\
    \bottomrule
  \end{tabular}
  \tablenote{注：表中粗体数值表示在当前对比评估项目下的最优性能指标。}
\end{table}

\section{性能指标对比与讨论}

由表~\ref{tab:exp_results}可以看出，相较于传统对比方案，本文算法在运算速度上提升了约 44\%，吞吐量提升了 39\% 以上。这一结果充分体现了自适应松弛机制在处理异构数据流时的明显优势。

\section{本章小结}

本章通过定性与定量对比实验，全面评估了新方法在不同场景下的综合表现，各项指标均达到了预期设计要求。
